Considereu el sistema d'equacions
\[
\left.\begin{aligned}
x+3y+az&=1\\
-2x-6y+2z&=2a\\
ax+3y+z&=3
\end{aligned}\right\},
\]
on $a$ és un paràmetre real, i siguin $\pi_1$ el pla determinat per la primera equació, $\pi_2$ el determinat
per la segona, i $\pi_3$ el determinat per la tercera.

\begin{apartats}

\apartat{1}
Discutiu el sistema en funció del valor de $a$.

\begin{solucio}
Comencem calculant el determinant de la matriu de coeficients del sistema:
\[
\begin{vmatrix}1&3&a\\-2&-6&2\\a&3&1\end{vmatrix}=-6+6a-6a-\left(-6a^2+6-6\right)=6a^2-6=6\left(a^2-1\right).
\]
Clarament, s'anul·la només quan $a=\pm1$. Per tant, si $a\neq1$ i $a\neq-1$, el rang de la matriu de coeficients
és 3, el rang de la matriu ampliada també i, com que tenim tres incògnites, és un sistema compatible
determinat.\\
Per al cas $a=1$, fem Gauss: a la fila 2 li sumem dues vegades la fila 1, i a la fila 3 li restem la fila 1:
\[
\left(\begin{array}{ccc|c}1&3&1&1\\-2&-6&2&2\\1&3&1&3\end{array}\right)\sim
\left(\begin{array}{ccc|c}1&3&1&1\\0&0&4&4\\0&0&0&2\end{array}\right).
\]
Clarament, veient la tercera equació, es tracta d'un sistema incompatible.\\
Finalment, per al cas $a=-1$, tornem a fer Gauss: a la fila 2 li sumem dues vegades la fila 1, i a la fila 3
li sumem la fila 1:
\[
\left(\begin{array}{ccc|c}1&3&-1&1\\-2&-6&2&-2\\-1&3&1&3\end{array}\right)\sim
\left(\begin{array}{ccc|c}1&3&-1&1\\0&0&0&0\\0&6&0&4\end{array}\right).
\]
Ara està clar que els rangs de les matrius de coeficients i ampliada són tots dos 2 i, com que hi ha tres
incògnites, el sistema és compatible indeterminat amb un grau de llibertat.

\textit{Pauta oficial:} 0,25 pel càlcul del determinant, i 0,25 per la discussió de cadascun dels tres casos.
\end{solucio}

\apartat{0,75}
Descriviu la posició relativa dels tres plans $\pi_1$, $\pi_2$ i $\pi_3$ en funció del valor de $a$.

\begin{solucio}
Per al cas $a\neq1$ i $a\neq-1$, el sistema ens ha donat compatible determinat, la qual cosa significa que els
tres plans $\pi_1$, $\pi_2$ i $\pi_3$ es tallen en un sol punt comú.\\
Per al cas $a=1$, el sistema és incompatible i, per tant, els tres plans alhora no s'intersequen. Clarament,
$\pi_1$ i $\pi_3$ són plans paral·lels (perquè els vectors $(1,3,1)$ i $(1,3,1)$ són proporcionals), i $\pi_2$
és un pla secant als dos anteriors.\\
Finalment, per al cas $a=-1$, tenim que $\pi_1=\pi_2$ (les equacions són proporcionals) i $\pi_3$ és un pla
secant a l'anterior; per tant, la intersecció dels tres és una recta.

\textit{Pauta oficial:} 0,25 per interpretar la posició relativa en cadascun dels tres casos.
\end{solucio}

\apartat{0,75}
En algun cas la intersecció dels tres plans és una recta? En cas afirmatiu, digueu per a quin valor del
paràmetre, i doneu un vector director i un punt d'aquesta recta.

\begin{solucio}
Pel que hem vist a l'apartat anterior, l'únic cas en què la intersecció dels tres plans és una recta és per a
$a=-1$. Resolem el sistema aprofitant els càlculs de l'apartat a). Tenim $y=\frac46=\frac23$ i
$x+3\cdot\frac23-z=1$; per tant, $x=z-1$. Les solucions, i per tant els punts de la recta en qüestió, són de
la forma
\[
(x,y,z)=\left(z-1,\tfrac23,z\right)=\left(-1,\tfrac23,0\right)+z(1,0,1).
\]
D'aquí veiem que $\left(-1,\frac23,0\right)$ és un punt d'aquesta recta i $(1,0,1)$ n'és un vector director.

\textit{Pauta oficial:} 0,25 per agafar el valor adequat del paràmetre, 0,25 per resoldre el sistema, i 0,25
per deduir un punt i un vector director.
\end{solucio}

\end{apartats}
